\documentclass[12pt]{article}

\usepackage{amsmath} %Advanced math typing
\usepackage{amsthm} %For theorem styles
\usepackage{amsfonts} %For \mathbb
\usepackage{amssymb} %For more math symbols
\usepackage{stmaryrd} %For even more math symbols (like \llbracket)
\usepackage{mathtools} %For extensible math symbols

\usepackage{array} %For special commands in tables
\usepackage[inline]{enumitem} %For enumerate styles

\usepackage[in]{fullpage} %Gets rid of stupidly large margins leaving 1 inch on every side

\usepackage[hypcap=true]{subcaption} % For captions in subfloats

\usepackage[%
  pdftex,% For direct pdf compilation compatibility
  plainpages=false,% Usage of not only arabic numbers
  pdfpagelabels,% Enable page labels
%  pagebackref,% Bibliographic back links to pages
  colorlinks,% Colored links
  citecolor=black,% Citation color
  linkcolor=black,% Link color
  urlcolor=black,% URL color
  filecolor=black,% File color
  bookmarksopen=false% For starting document with bookmarks tree closed
]{hyperref} %For pdf links
\usepackage[all]{hypcap} %Fixes figures and tables anchors for hyperref

% Theorem styles
%----------------------------------------------------------------
\newtheoremstyle{thmstyle}%style name
  {\medskipamount}%space before
  {\smallskipamount}%space after
  {\slshape}%font used
  {0pt}%indentation
  {\bfseries}%modifier theorem head
  {.}%punctuation between theorem head and body
  { }%space after punctuation
  {\thmname{#1}\thmnumber{ #2}{\normalfont\thmnote{ (#3)}}}%theorem specifier

\newtheoremstyle{bigskipthmstyle}%style name
  {\medskipamount}%space before
  {\bigskipamount}%space after
  {\slshape}%font used
  {0pt}%indentation
  {\bfseries}%modifier theorem head
  {.}%punctuation between theorem head and body
  { }%space after punctuation
  {\thmname{#1}\thmnumber{ #2}{\normalfont\thmnote{ (#3)}}}%theorem specifier
  
\newtheoremstyle{plainstyle}%style name
  {\medskipamount}%space before
  {\smallskipamount}%space after
  {\rmfamily}%font used
  {0pt}%indentation
  {\bfseries}%modifier theorem head
  {.}%punctuation between theorem head and body
  { }%space after punctuation
  {\thmname{#1}\thmnumber{ #2}{\normalfont\thmnote{ (#3)}}}%theorem specifier

\theoremstyle{bigskipthmstyle}
\newtheorem{topic}{Topic}

\theoremstyle{thmstyle}
\newtheorem{theorem}{Theorem}[section]
\newtheorem{lemma}[theorem]{Lemma}
\newtheorem{proposition}[theorem]{Proposition}
\newtheorem{claim}{Claim}[topic]
\newtheorem{fact}[theorem]{Fact}
\newtheorem{axiom}[theorem]{Axiom}


\theoremstyle{plainstyle}
\newtheorem{definition}[theorem]{Definition}
%----------------------------------------------------------------

\newenvironment{proofof}[1]{\begin{proof}[Proof of #1.]}{\end{proof}}

\newcommand{\floor}[1]{\lfloor #1\rfloor}
\newcommand{\ceil}[1]{\lceil #1\rceil}

\setlist[enumerate]{label={\roman*.}, ref={(\roman*)}}

\AddToHook{cmd/part/before}{\clearpage} % Forces each part to be in a new page
\AddToHook{cmd/section/before}{\clearpage} % Forces each section to be in a new page

% General symbols
\newcommand{\df}{\stackrel{\text{def}}{=}}
\newcommand{\rest}{\mathord{\vert}}
\newcommand{\divides}{\mathbin{\vert}}

% Common operators
\DeclareMathOperator{\LP}{LP}
\DeclareMathOperator{\dom}{dom}

% Math-bold symbols
\newcommand{\FF}{\mathbb{F}}
\newcommand{\NN}{\mathbb{N}}
\newcommand{\QQ}{\mathbb{Q}}
\newcommand{\RR}{\mathbb{R}}
\newcommand{\ZZ}{\mathbb{Z}}

% Calligraphic symbols
\newcommand{\cB}{\mathcal{B}}
\newcommand{\cC}{\mathcal{C}}
\newcommand{\cF}{\mathcal{F}}
\newcommand{\cG}{\mathcal{G}}
\newcommand{\cH}{\mathcal{H}}
\newcommand{\cK}{\mathcal{K}}
\newcommand{\cP}{\mathcal{P}}

% Names
\def\Schroder{Schr\"{o}der}
\def\Bezout{B\'{e}zout}



\begin{document}

\section{Tychonoff's Theorem}

The objective of this advanced topic is to introduce you to enough concepts of topology to be able
to prove Tychonoff's Theorem on compactness.

\subsection{Preliminaries}

To prove Tychonoff's Theorem, we will need the Axiom of Choice, so we start by setting up some
notation and reminding the axiom.
\begin{definition}
  If $I$ is a set and $(X_i)_{i\in I}$ is a collection of sets indexed by $I$, then $\prod_{i\in I}
  X_i$ is the set of all functions $f\colon I\to\bigcup_{i\in I} X_i$ such that for every $i\in I$,
  we have $f(i)\in X_i$.

  We typically think of elements of $\prod_{i\in I} X_i$ as sequences indexed by $I$ whose $i$th
  element is in $X_i$, that is, $f\in\prod_{i\in I} X_i$ corresponds to the sequence
  $x\in\prod_{i\in I} X_i$ given by $x_i\df f(i)$.

  Furthermore, with a small abuse of notation, if for every $i\in I$, we have $Y_i\subseteq X_i$,
  then we think of $\prod_{i\in I} Y_i$ as a subset of $\prod_{i\in I} X_i$ by simply noting that if
  $f\in\prod_{i\in I} Y_i$, then we can simply extend the codomain of $f$ to $\bigcup_{i\in I} X_i$
  to think of $f$ as an element of $\prod_{i\in I} X_i$; in other words, $\prod_{i\in I} Y_i$ is the
  subset of all $f\in\prod_{i\in I} X_i$ such that for every $i\in I$, we have $f(i)\in Y_i$.
\end{definition}

The Axiom of Choice below might seem obviously true but it does \emph{not} follow from basic
set-theoretic axioms (the non-triviality of the Axiom of Choice is precisely when $I$ is infinite):
\begin{axiom}[Axiom of Choice]
  If $I$ is a set and $(X_i)_{i\in I}$ is a collection of sets indexed by $I$ such that every $X_i$
  is non-empty, then $\prod_{i\in I} X_i$ is non-empty.
\end{axiom}

We will also need Zorn's Lemma, which is equivalent to the Axiom of Choice, so we recall the
necessary definitions below.

\begin{definition}
  A \emph{partial order} on a set $X$ is a relation ${\lhd}\subseteq X\times X$ satisfying the
  following (as usual, we make the small abuse of notation and write ``$x\lhd y$'' to mean
  ``$(x,y)\in{\lhd}$''):
  \begin{description}
  \item[Reflexivity:] For every $x\in X$, we have $x\lhd x$.
  \item[Antisymmetry:] For every $x,y\in X$, if $x\lhd y$ and $y\lhd x$, then $x = y$.
  \item[Transitivity:] For every $x,y,z\in X$, if $x\lhd y$ and $y\lhd z$, then $x\lhd z$.
  \end{description}

  If $\lhd$ is a partial order on $X$ and $m\in X$, we say that $m$ is \emph{maximal} for $\lhd$ if
  for every $x\in X$, if $m\lhd x$, then $m=x$.

  If $\lhd$ is a partial order on $X$ and $A\subseteq X$, then the restriction of $\lhd$ to $A$ is
  the set ${\lhd\rest_A}\df{\lhd}\cap A\times A$ (i.e., we simply ignore what $\lhd$ does on
  points outside of $A$).

  We say that a partial order $\lhd$ on $X$ is a \emph{linear order} if it further satisfies:
  \begin{description}
  \item[Totality:] For every $x,y\in X$, we have $x\lhd y$ or $y\lhd x$.
  \end{description}

  We say that a partial order $\lhd$ on $X$ satisfies \emph{Zorn's condition} if for every set
  $A\subseteq X$, if $\lhd\rest_A$ is a linear order on $A$, then there exists $u\in X$ such that
  for every $a\in A$, we have $a\lhd u$. (We abuse a bit the terminology in the above by saying that
  $A$ is linearly ordered instead of saying that $\lhd\rest_A$ is a linear order on $A$.)
\end{definition}

\begin{lemma}[Zorn's Lemma]\label{lem:Zorn}
  Let $\lhd$ be a partial order on $X$ satisfying Zorn's condition. Then (assuming the Axiom of
  Choice holds) there exists $m\in X$ such that $m$ is maximal for $\lhd$.  
\end{lemma}

We start with the definition of topologies and topological spaces (which you might remember from one
of the exams in IBL~1) and the abstract definition of compactness:
\begin{definition}
  A \emph{topology} on a set $X$ is a collection $\tau\subseteq\cP(X)$ of subsets of $X$ satisfying
  the following properties:
  \begin{description}
  \item[Trivial open sets:] $\varnothing$ and $X$ are elements of $\tau$.
  \item[Closure under unions:] If $(G_i)_{i\in I}$ is a collection of sets in $\tau$, then the union
    $\bigcup_{i\in I} G_i$ is in $\tau$.
  \item[Closure under finite intersections:] If $G_1,G_2\in\tau$, then $G_1\cap G_2\in\tau$.
  \end{description}

  Sets in $\tau$ are called \emph{open sets} (with respect to the topology $\tau$). A set
  $F\subseteq X$ whose complement $X\setminus F$ is open is called \emph{closed} (with respect to
  $\tau$).

  A \emph{topological space} is a pair $(X,\tau)$, where $\tau$ is a topology on $X$.

  Given a topological space $(X,\tau)$, as usual, an open cover of a set $Y\subseteq X$ is a
  collection $\cG\subseteq\tau$ (i.e., a collection of open sets) such that $\bigcup_{G\in\cG}
  G\supseteq Y$. A \emph{subcover} of $\cG$ is a set $\cG'\subseteq\cG$ that is also a cover of $Y$.

  As usual, set $Y\subseteq X$ is \emph{compact} (with respect to $\tau$) if every open cover of $Y$
  has a finite subcover.
\end{definition}

The two most simple topologies on a space $X$ are the \emph{discrete topology} $\cP(X)$ (in which
every set is open) and the \emph{trivial topology} $\{\varnothing, X\}$ (in which only $\varnothing$
and $X$ are open).

Tychonoff's Theorem can be summarized by the sentence: products of compact sets are compact in the
product topology. We define the product topology below.

\begin{definition}
  Let $I$ be a set and suppose $(X_i,\tau_i)_{i\in I}$ is a collection of topological spaces. A
  \emph{cylinder open set} of $\prod_{i\in I} X_i$ is a subset of the form $\prod_{i\in I} U_i$ such
  that each $U_i\in\tau_i$ is an open set of $X_i$ and
  \begin{align*}
    \{i\in I \mid U_i\neq X_i\}
  \end{align*}
  is finite (that is, all except for finitely many $U_i$ are equal to $X_i$).

  A \emph{cylinder closed set} is defined analogously, but we require the $U_i$ to be closed sets.

  A \emph{basic closed set} is any set that can be written as a finite union of cylinder closed sets.

  The \emph{product topology} is the collection of all subsets of $\prod_{i\in I} X_i$ that can be
  written as (finite or infinite) unions of cylinder open sets.
\end{definition}

\begin{lemma}
  Cylinder closed sets and basic closed sets are closed sets (in the product topology).
\end{lemma}

\begin{lemma}\label{lem:compopen}
  The complement of a cylinder open set is a basic closed set (in the product topology).
\end{lemma}

\begin{lemma}\label{lem:closed}
  A set is closed in the product topology if and only if it can be written as a (finite or infinite)
  intersection of basic closed sets.
\end{lemma}

\subsection{Main topic}

First, we will see an alternative formulation of compactness via intersection properties.

\begin{definition}
  A collection of sets $\cF$ is said to have \emph{finite intersection property} if for every
  \emph{finite} non-empty subcollection $\cF'\subseteq\cF$, the intersection $\bigcap_{F\in\cF'} F$
  is non-empty (i.e., every finite non-empty subcollection of $\cF$ has a common point).

  We say that $\cF$ has \emph{finite intersection property within $Y$} if $\cF\cup\{Y\}$ has finite
  intersection property (note that this is equivalent to saying that every finite non-empty
  subcollection of $\cF$ has a common point within $Y$).
\end{definition}

\begin{theorem}\label{thm:FIP}
  Let $(X,\tau)$ be a topological space. The following are equivalent for set $Y\subseteq X$:
  \begin{enumerate}
  \item $Y$ is compact (with respect to $\tau$).
  \item For every collection $\cF$ of closed sets with finite intersection property within $Y$, the
    intersection $Y\cap\bigcap_{F\in\cF} F$ is non-empty.
  \end{enumerate}
  
\begin{proof}
\vspace{0.5em}
\textbf{(i) $\Rightarrow$ (ii):} Assume $Y$ is compact. Let 
\(\mathcal{F}\) be a collection of closed sets with the finite intersection property 
within \(Y\). We wish to show that
\[
Y \cap \bigcap_{F\in \mathcal{F}} F \neq \varnothing.
\]

Suppose, for the sake of contradiction, that
\[
Y \cap \bigcap_{F\in \mathcal{F}} F = \varnothing.
\]
Then, for every \(y\in Y\) there exists some \(F\in \mathcal{F}\) such that 
\(y\notin F\). Equivalently, \(y\in X\setminus F\). Since each \(F\) is closed, 
\(X\setminus F\) is open. Hence, 
\[
\{\, X\setminus F \mid F\in \mathcal{F} \,\}
\]
forms an open cover of \(Y\).

By the compactness of \(Y\), there exists a finite subcover; that is, there exist 
\(F_1, F_2, \dots, F_n \in \mathcal{F}\) such that
\[
Y \subseteq \bigcup_{i=1}^n \Bigl(X \setminus F_i\Bigr).
\]
By De Morgan's Laws, we have
\[
\bigcup_{i=1}^n \Bigl(X \setminus F_i\Bigr)
=
X\setminus \Bigl(\bigcap_{i=1}^n F_i\Bigr).
\]
Thus, 
\[
Y \subseteq X \setminus \Bigl(\bigcap_{i=1}^n F_i\Bigr),
\]
which implies that
\[
Y \cap \Bigl(\bigcap_{i=1}^n F_i\Bigr) = \varnothing.
\]
This contradicts the assumption that \(\mathcal{F}\) has the finite intersection property within \(Y\). 

\vspace{0.5em}
\textbf{(ii) $\Rightarrow$ (i):} Now, assume that for every collection 
\(\mathcal{F}\) of closed sets with the FIP within \(Y\) we have
\[
Y \cap \bigcap_{F\in \mathcal{F}} F \neq \varnothing.
\]
We want to show that \(Y\) is compact.

Suppose, toward a contradiction, that \(Y\) is not compact. Then there exists an open 
cover \(\mathcal{G}=\{G_i\}_{i\in I}\) of \(Y\) which does not admit any finite subcover.

Since each \(G_i\) is open, the complement \(X\setminus G_i\) is closed. Define
\[
\mathcal{H} = \{\,X\setminus G_i \mid i\in I\,\}.
\]
We claim that \(\mathcal{H}\) has the finite intersection property within \(Y\).

Indeed, take any finite subcollection \(\{X\setminus G_{i_1}, X\setminus G_{i_2}, \dots, X\setminus G_{i_n}\}\) of \(\mathcal{H}\). Then,
\[
\bigcap_{j=1}^n \Bigl(X\setminus G_{i_j}\Bigr)
=
X\setminus \Bigl(\bigcup_{j=1}^n G_{i_j}\Bigr).
\]
Since \(\{G_{i_j}\}_{j=1}^n\) is a finite subset of the open cover \(\mathcal{G}\), and by assumption no finite subcollection of \(\mathcal{G}\) covers \(Y\), it follows that
\[
Y \setminus \Bigl(\bigcup_{j=1}^n G_{i_j}\Bigr) \neq \varnothing.
\]
Hence,
\[
Y \cap \Bigl(\bigcap_{j=1}^n (X\setminus G_{i_j})\Bigr) \neq \varnothing.
\]
Thus, \(\mathcal{H}\) has the FIP within \(Y\).

By our hypothesis (applied to the collection \(\mathcal{H}\) of closed sets), we have
\[
Y \cap \bigcap_{H\in \mathcal{H}} H \neq \varnothing.
\]
That is, there exists a point
\[
m \in Y \cap \bigcap_{i\in I} (X\setminus G_i).
\]
In other words, \(m\in Y\) and for all \(i\in I\),
\[
m \notin G_i.
\]
This contradicts the fact that \(\mathcal{G}\) is an open cover of \(Y\).

Thus, \(Y\) must be compact.


\renewcommand\qedsymbol{QED}
\end{proof}


\end{theorem}

\begin{lemma}\label{lem:bcFIP}
  Let $(X_i,\tau_i)_{i\in I}$ be a collection of topological spaces and let $Y\subseteq\prod_{i\in I} X_i$. Suppose every
  collection $\cF$ of basic closed sets with finite intersection property within $Y$ has non-empty intersection.

  Then $Y$ is compact.

  
\begin{proof}
By Theorem 1.11 it suffices to show that every collection $\mathcal{F}$ of closed sets with the FIP within $Y$ satisfies
\[
Y \cap \bigcap_{F\in\mathcal{F}} F \neq \varnothing.
\]

Let $\mathcal{F}$ be any collection of closed sets in $\prod_{i\in I} X_i$ that has the FIP within $Y$. By Lemma 1.9, every closed set $F\in\mathcal{F}$ can be expressed as an (arbitrary) intersection of basic closed sets. That is, for each $F\in\mathcal{F}$ there exists an index set $J_F$ and basic closed sets $\{G^F_j : j\in J_F\}$ such that
\[
F = \bigcap_{j\in J_F} G^F_j.
\]
Now, define the collection
\[
\mathcal{G} = \bigcup_{F\in\mathcal{F}} \{G^F_j : j\in J_F\}.
\]
We claim that $\mathcal{G}$ has the FIP within $Y$.

To see this, let $\{G_1, G_2, \dots, G_n\}$ be an arbitrary finite subcollection of $\mathcal{G}$. For each $k=1,\dots,n$, there exists some $F_k\in\mathcal{F}$ and an index $j_k\in J_{F_k}$ with 
\[
G_k = G^{F_k}_{j_k}.
\]
Then,
\[
\bigcap_{k=1}^n F_k \subseteq \bigcap_{k=1}^n G^{F_k}_{j_k} = \bigcap_{k=1}^n G_k.
\]
Since $\mathcal{F}$ has the FIP within $Y$, we have
\[
Y\cap \bigcap_{k=1}^n F_k \neq \varnothing.
\]
Thus, it follows that
\[
Y\cap \bigcap_{k=1}^n G_k \neq \varnothing.
\]
This verifies that every finite subcollection of $\mathcal{G}$ has nonempty intersection with $Y$.

By hypothesis, since $\mathcal{G}$ is a collection of basic closed sets with the FIP within $Y$, we conclude that
\[
Y \cap \bigcap_{G\in \mathcal{G}} G \neq \varnothing.
\]
Let $x$ be any point in $Y \cap \bigcap_{G\in \mathcal{G}} G$. Now, fix an arbitrary $F\in \mathcal{F}$. By construction, 
\[
F = \bigcap_{j\in J_F} G^F_j,
\]
and each $G^F_j$ is an element of $\mathcal{G}$. Hence, $x$ lies in every $G^F_j$ and, consequently, in
\[
F = \bigcap_{j\in J_F} G^F_j.
\]
Since $F$ was arbitrary in $\mathcal{F}$, we deduce that
\[
x\in \bigcap_{F\in \mathcal{F}} F.
\]
It follows that
\[
Y\cap \bigcap_{F\in \mathcal{F}} F \neq \varnothing.
\]

By Theorem 1.11, this shows that $Y$ is compact.

\renewcommand\qedsymbol{QED}
\end{proof}


\end{lemma}

\begin{lemma}\label{lem:cylFIP}
  Assume the Axiom of Choice holds. Let $(X_i,\tau_i)_{i\in I}$ be a collection of topological
  spaces and for each $i\in I$, let $Y_i\subseteq X_i$ be a compact set (with respect to $\tau_i$).

  Then every collection $\cF$ of cylinder closed sets of $\prod_{i\in I} X_i$ with finite intersection property within
  $\prod_{i\in I} Y_i$ has non-empty intersection within $\prod_{i\in I} Y_i$.

  Hint: Write each $F\in\cF$ as $\prod_{i\in I} K^F_i$ with $K^F_i\subseteq X_i$ closed and for each $i\in I$, show that the
  family $\cF_i\df\{K^F_i \mid F\in\cF\}$ has finite intersection property within $Y_i$. Use compactness of $Y_i$ to
  $Y_i\cap\bigcap_{F\in\cF_i}$ and use the Axiom of Choice to construct an element in $Y\cap\bigcap_{F\in\cF} F$.

  
\begin{proof}
Let $\cF$ be a collection of cylinder closed sets with FIP in 
\[
\prod_{i\in I} Y_i.
\]
For any $F\in \cF$, we can write
\[
F = \prod_{i\in I} K^F_i,
\]
with $K^F_i \subseteq X_i$ closed.

For each fixed $i\in I$, define
\[
\cF_i \df \{K^F_i \mid F\in \cF\}.
\]
Take a finite subcollection $\cF'_i\subseteq \cF_i$ and write
\[
\cF'_i = \{K^{F_1}_i, K^{F_2}_i, \dots, K^{F_n}_i\}.
\]
Since each $K^{F_m}_i$ comes from some $F_m\in \cF$, let
\[
\cG = \{F_1, F_2, \dots, F_n\},
\]
so that $\cG$ is a finite subcollection of $\cF$.

By the FIP of $\cF$, we have
\[
\left( \bigcap_{G \in \cG} G \right) \cap \left( \prod_{i\in I} Y_i \right) \neq \varnothing.
\]
Thus, there exists an element 
\[
a \in \prod_{i\in I} Y_i
\]
such that
\[
a \in \bigcap_{G \in \cG} G.
\]
In particular, for each $i\in I$, if we denote the $i$th coordinate of $a$ by $a_i$, then
\[
a_i \in Y_i \quad \text{and} \quad a_i \in \bigcap_{m=1}^n K^{F_m}_i.
\]
This shows that $\cF_i$ has the finite intersection property within $Y_i$.

Since each $Y_i$ is compact, it follows that for every $i\in I$,
\[
Y_i \cap \bigcap_{F\in \cF} K^F_i \neq \varnothing.
\]
Thus, for each $i\in I$, there exists an element 
\[
h_i \in Y_i \cap \bigcap_{F\in \cF} K^F_i.
\]

By the Axiom of Choice, we can select such an $h_i$ for every $i\in I$ and form
\[
h = (h_i)_{i\in I} \in \prod_{i\in I} Y_i.
\]
Since for each $F\in \cF$, we have $h_i \in K^F_i$ for all $i\in I$, it follows that
\[
h \in \bigcap_{F\in \cF} F.
\]
Thus,
\[
h \in \left( \prod_{i\in I} Y_i \right) \cap \left( \bigcap_{F\in \cF} F \right),
\]
showing that the intersection is non-empty.

\renewcommand\qedsymbol{QED}
\end{proof}


\end{lemma}

\begin{lemma}\label{lem:bccyl}
  Assume the Axiom of Choice holds. Let $(X_i,\tau_i)_{i\in I}$ be a collection of topological
  spaces and let $Y\subseteq\prod_{i\in I} X_i$.

  Suppose $\cF$ is a collection of basic closed sets of $\prod_{i\in I} X_i$ with finite
  intersection property within $Y$. Then there exists a function $h$ that maps each element $F$ of
  $\cF$ to a cylinder closed set $h(F)\subseteq F$ such that $h(\cF)$ has finite intersection
  property within $Y$.

  Hint: Let $\cC$ be the set of all cylinder closed sets of $\prod_{i\in I} X_i$ and let $\cG$ be the set of all functions of
  the form $g\colon\cK\to\cC$ with $\cK\subseteq\cF$ such that $g(F)\subseteq F$ for every $F\in\cF$ and
  \begin{align*}
    \cF_g & \df g(\cK)\cup(\cF\setminus\cK)
  \end{align*}
  has finite intersection property within $Y$. Equip $\cG$ with the partial order of function extension $\subseteq$, i.e.,
  $g_1\subseteq g_2$ if $\dom(g_1)\subseteq\dom(g_2)$ and $g_2\rest_{\dom(g_1)} = g_1$. Prove that $\cG$ satisfies Zorn's
  condition, apply Zorn's Lemma to get a maximal function $h$ and argue by contradiction that we must have $\dom(h) = \cF$, so
  $h$ is our desired function.


\end{lemma}

\begin{theorem}[Tychonoff's Theorem]
  Assume the Axiom of Choice holds. Let $(X_i,\tau_i)_{i\in I}$ be a collection of topological
  spaces and for each $i\in I$, let $Y_i\subseteq X_i$ be a compact set (with respect to $\tau_i$).

  Then $\prod_{i\in I} Y_i$ is a compact set in $\prod_{i\in I} X_i$ (with respect to the product
  topology).
  
  
\begin{proof}[Short Proof]
By Lemma~1.12, it suffices to show that every collection of \emph{basic closed} sets in 
\(\prod_{i\in I} X_i\) with the finite intersection property (FIP) within 
\(\prod_{i\in I} Y_i\) has a nonempty intersection.

By Lemma~1.14, given any such collection of basic closed sets with the FIP, we can assign 
to each one a \emph{cylinder closed} subset so that the new family still has the FIP in 
\(\prod_{i\in I} Y_i\).

Then, Lemma~1.13 tells us that any collection of cylinder closed sets with the FIP in 
\(\prod_{i\in I} Y_i\) has nonempty intersection. Hence every family of basic closed sets 
with the FIP also has nonempty intersection, and by Lemma~1.12, \(\prod_{i\in I} Y_i\) is compact.

\renewcommand\qedsymbol{QED}
\end{proof}


\end{theorem}

A direct consequence of Tychonoff's Theorem is the Compactness Theorem for Propositional Logic,
which says that if one has a system with infinitely many true/false variables and infinitely many
constraints (each involving finitely many variables) such that every finite subsystem has a
solution, then the full system has a solution. This essentially amounts to noticing that $\{0,1\}$
equipped with the discrete topology is compact, so an infinite product $\prod_{i\in I}\{0,1\}$ is
also compact, any constraint amounts to a cylinder closed set and finite subsystems having a
solution amounts to finite intersection property.

\end{document}

%% Local Variables:
%% mode: latex
%% tab-width: 2
%% End:
